% Part: first-order-logic
% Chapter: models-theories
% Section: set-theory

\documentclass[../../../include/open-logic-section]{subfiles}

\begin{document}

\olfileid{fol}{mat}{set}

\olsection{సమితుల సిద్ధాంతం}

దాదాపు మొత్తం గణితాన్ని సమితుల సిద్ధాంతంలో అభివృద్ధి చేయవచ్చు. ఈ
సిద్ధాంతంలో గణితాన్ని అభివృద్ధి చేయడానికి అనేక పనులు అవసరం. మొదటగా,
$\in$ సంబంధానికి స్వీకృతాల సమితి అవసరం. కొన్నిసార్లు~$\in$కు పరస్పర
విరుద్ధమైన ధర్మాలను ఇచ్చే అనేక స్వీకృత వ్యవస్థలు అభివృద్ధి చేయబడ్డాయి.
ఎంపిక స్వీకృతంతో కూడిన జెర్మెలో--ఫ్రెంకెల్ సమితి సిద్ధాంతమైన
$\Log{ZFC}$ అనే స్వీకృత వ్యవస్థ వాటిలో ప్రత్యేకంగా నిలుస్తుంది: గణితవేత్తలు
నిరూపించగలమని ఆశించే దాదాపు అన్నిటినీ నిరూపించడానికి దాని స్వీకృతాలు
సరిపోతాయని తేలడం వల్ల అది అత్యంత విస్తృతంగా ఉపయోగించబడిన, అధ్యయనం
చేయబడిన వ్యవస్థ. కానీ దీనిని స్థాపించడానికి ముందు, గణితవేత్తలు వ్యక్తం
చేయదలచిన విషయాలన్నిటినీ అసలు ఎలా \emph{వ్యక్తం} చేయగలమో స్పష్టం
చేయాలి. మొదటగా, ఈ భాషలో \tetoken{స్థిరసంకేతాలు}{!!{constant}s} లేదా \tetoken{ప్రమేయ సంకేతాలు}{!!{function}s} లేవు. కనుక
నిర్దిష్ట సమితుల (ఉదాహరణకు $\emptyset$ లేదా $\Nat$) గురించి, సమితులపై
చర్యల (ఉదాహరణకు $X \cup Y$, $\Pow{X}$) గురించి ఎలా మాట్లాడగలమో తొలి
చూపులో స్పష్టంగా ఉండదు; సంబంధాలు, ప్రమేయాలు వంటి సమితులు కాని వాటిని
కలిగిన నిర్మాణాల సంగతి చెప్పనే అక్కర్లేదు.

``దీనిలో మూలకం'' అనేది మాత్రమే మనకు ఆసక్తి ఉన్న సంబంధం కాదు:
``దీనికి ఉపసమితి'' కూడా దాదాపు అంతే ముఖ్యమైనదిగా కనిపిస్తుంది. కానీ
``దీనికి ఉపసమితి''ని ``దీనిలో మూలకం'' ద్వారా \emph{నిర్వచించవచ్చు}.
దీనికోసం, సమితుల సిద్ధాంత భాషలోని \tetoken{ఒక సూత్రం}{!!a{formula}}~$!A(x,y)$ను
కనుగొనాలి; సమితుల జత~$\tuple{X,Y}$, $X \subseteq Y$ అయితేనే
దానిని సంతృప్తిపరచాలి. $X$లోని \tetoken{మూలకాలు}{!!{element}s} అన్నీ $Y$లో కూడా
\tetoken{మూలకాలు}{!!{element}s} అయితేనే $X$, $Y$కు ఉపసమితి. కాబట్టి
$\subseteq$ను కింది సూత్రం ద్వారా నిర్వచించవచ్చు:
\[
\lforall[z][(z \in x \lif z \in y)]
\]
ఇక నుంచి ఒక సూత్రంలో~$\subseteq$ సంబంధాన్ని ఉపయోగించదలచినప్పుడల్లా,
దాని స్థానంలో ఈ సూత్రాన్ని ఉపయోగించవచ్చు ($x$, $y$లను తగిన విధంగా
మార్చి, అవసరమైతే బద్ధ చరం~$z$కు మరో పేరు పెట్టాలి). ఉదాహరణకు, సమితుల
విస్తారత అంటే ఏ సమితులు~$x$, $y$ అయినా ఒకటి మరొకదానిలోను, మరొకటి
మొదటిదానిలోను ఇమిడి ఉంటే, $x$, $y$ ఒకే సమితి కావాలి. దీనిని
$\lforall[x][\lforall[y][((x \subseteq y \land y
    \subseteq x) \lif x = y)]]$ ద్వారా వ్యక్తం చేయవచ్చు; లేదా పై
నిర్వచనంతో~$\subseteq$ను ప్రతిస్థాపిస్తే, కింది విధంగా వ్యక్తం చేయవచ్చు:
\[
\lforall[x][\lforall[y][((\lforall[z][(z \in x \lif z \in y)] \land
    \lforall[z][(z \in y \lif z \in x)]) \lif x = y)]].
\]
నిజానికి ఇది~$\Log{ZFC}$ స్వీకృతాల్లో ఒకటైన ``విస్తారతా స్వీకృతం.''

$\emptyset$కు \tetoken{స్థిరసంకేతం}{!!{constant}} లేదు; కానీ ``$x$ శూన్యం'' అని
$\lnot \lexists[y][y \in x]$ ద్వారా వ్యక్తం చేయవచ్చు. అప్పుడు
``$\emptyset$ ఉనికిలో ఉంది'' అనేది \tetoken{వాక్యం}{!!{sentence}}
$\lexists[x][\lnot \lexists[y][y \in x]]$ అవుతుంది. ఇది~$\Log{ZFC}$లో
మరో స్వీకృతం. (ఒకే ఒక్క శూన్య సమితి ఉందని విస్తారతా స్వీకృతం
సూచిస్తుందని గమనించండి.) సమితుల సిద్ధాంత భాషలో~$\emptyset$ గురించి
మాట్లాడదలచినప్పుడల్లా, ``శూన్యమైన ఒక సమితి ఉంది, ఇంకా~\dots'' అని
రాస్తాం. ఉదాహరణకు, $\emptyset$ ప్రతి సమితికీ ఉపసమితి అనే వాస్తవాన్ని
వ్యక్తం చేయడానికి ఇలా రాయవచ్చు:
\[
\lexists[x][(\lnot\lexists[y][y \in x] \land \lforall[z][x \subseteq
    z])]
\]
ఇక్కడ $x \subseteq z$ను కూడా దాని నిర్వచనంతో ప్రతిస్థాపించవలసిందే.

$X \cup Y$, $\Pow{X}$ వంటి సమితులపై చర్యల గురించి మాట్లాడడానికి కూడా
ఇలాంటి ఉపాయమే ఉపయోగించాలి. సమితుల సిద్ధాంత భాషలో ప్రమేయ సంకేతాలు లేవు;
కానీ $X \cup Y = Z$, $\Pow{X} = Y$ అనే ప్రమేయాత్మక సంబంధాలను కింది
వాటి ద్వారా వ్యక్తం చేయవచ్చు:
\begin{align*}
& \lforall[u][((u \in x \lor u \in y) \liff u \in z)]\\
& \lforall[u][(u \subseteq x \liff u \in y )]
\end{align*}
ఎందుకంటే $X \cup Y$లోని \tetoken{మూలకాలు}{!!{element}s} సరిగ్గా $X$లో
\tetoken{మూలకాలుగా}{!!{element}s} లేదా $Y$లో \tetoken{మూలకాలుగా}{!!{element}s} ఉన్న సమితులే;
$\Pow{X}$లోని \tetoken{మూలకాలు}{!!{element}s} సరిగ్గా
$X$ యొక్క ఉపసమితులే. అయితే దీనివల్ల $x \cup y$ లేదా $\Pow{x}$ను
పదాల్లా ఉపయోగించలేం: $X \cup Y = Z$, $\Pow{X} = Y$ సంబంధాలను
నిర్వచించే పూర్తి \tetoken{సూత్రాలను}{!!{formula}s} మాత్రమే ఉపయోగించగలం. నిజానికి ఈ
సంబంధాలు ఎప్పుడైనా సంతృప్తి చెందుతాయో, అంటే సమ్మేళనాలు, ఘాత సమితులు
ఎల్లప్పుడూ ఉంటాయో, మనకు ఇంకా తెలియదు. ఉదాహరణకు, \tetoken{వాక్యం}{!!{sentence}}
$\lforall[x][\lexists[y][\Pow{x} = y]]$, $\Log{ZFC}$లోని మరో
స్వీకృతం (ఘాత సమితి స్వీకృతం).

ఇప్పుడు క్రమయుగ్మాలు లేదా ప్రమేయాల గురించి ఎలా మాట్లాడాలి? ఇక్కడ
క్రమయుగ్మాలు, ప్రమేయాలను ప్రత్యేక రకాల సమితులుగా ఎలా భావించవచ్చో
వివరించాలి. క్రమయుగ్మం~$\tuple{x,y}$ను
$\{\{x\},\{x,y\}\}$ అనే సమితిగా నిర్వచించడం ఒక మార్గం. కానీ ఇంతకు
ముందులాగే, ఈ సమితికి పేరు పెట్టే \tetoken{ఒక ప్రమేయ సంకేతాన్ని}{!!a{function}} ప్రవేశపెట్టలేం;
$\tuple{x,y}=z$, అంటే $\{\{x\},\{x,y\}\}=z$ అనే సంబంధాన్ని మాత్రమే
నిర్వచించగలం:
\[
\lforall[u][(u \in z \liff (\lforall[v][(v \in u \liff v = x)] \lor
  \lforall[v][(v \in u \liff (v = x \lor v = y))]))]
\]
ఇది~$z$లోని \tetoken{మూలకాలు}{!!{element}s}~$u$, సరిగ్గా $x$ మాత్రమే \tetoken{మూలకంగా}{!!{element}}
ఉన్న లేదా $x$, $y$ మాత్రమే \tetoken{మూలకాలుగా}{!!{element}s} ఉన్న సమితులే అని చెబుతుంది
(మరోలా చెప్పాలంటే, $\{x\}$కు లేదా $\{x,y\}$కు సమానమైన సమితులే).
ఇది లభించిన తరువాత, ఉదాహరణకు $X \times Y = Z$ అని ఇంకా చెప్పవచ్చు:
\[
\lforall[z][(z \in Z \liff \lexists[x][\lexists[y][(x \in
        X \land y \in Y \land \tuple{x, y} = z)]])]
\]

ప్రమేయం~$f \colon X \to Y$ను $f(x)=y$ అనే సంబంధంగా, అంటే
$\Setabs{\tuple{x,y}}{f(x)=y}$ అనే యుగ్మాల సమితిగా భావించవచ్చు. అప్పుడు
సమితి~$f$, $X$ నుంచి $Y$కు ప్రమేయం అని చెప్పడానికి: (a) అది ఒక
సంబంధం~$\subseteq X \times Y$ కావాలి; (b) అది సర్వత్ర నిర్వచితమైనది
కావాలి, అంటే ప్రతి $x \in X$కు $\tuple{x,y}\in f$ అయ్యే ఏదో
$y \in Y$ ఉండాలి; (c) అది ప్రమేయాత్మకమైనది కావాలి, అంటే
$\tuple{x,y},\tuple{x,y'}\in f$ అయినప్పుడల్లా $y=y'$ కావాలి (ఎందుకంటే
ప్రమేయ విలువలు ఏకైకంగా ఉండాలి). కాబట్టి ``$f$, $X$ నుంచి $Y$కు ఒక
ప్రమేయం'' అని ఇలా రాయవచ్చు:
\begin{align*}
 \lforall[u][(u \in f \lif {}] & \lexists[x][\lexists[y][(x \in X \land y \in
      Y \land \tuple{x, y} = u)]]) \land {}\\
 \lforall[x][(x \in X \lif {}] &
      (\lexists[y][(y \in Y \land \mathrm{maps}(f, x, y))] \land {}\\
& (\lforall[y][\lforall[y'][((\mathrm{maps}(f, x, y) \land
    \mathrm{maps}(f, x, y')) \lif y = y')]]))
\end{align*}
ఇక్కడ $\mathrm{maps}(f,x,y)$ అనేది
$\lexists[v][(v \in f \land \tuple{x,y}=v)]$కు సంక్షిప్త రూపం (ఈ
\tetoken{సూత్రం}{!!{formula}}, ``$f(x)=y$'' అని వ్యక్తం చేస్తుంది).

ఉదాహరణకు, $f\colon X\to Y$ \tetoken{ఒకటి-ఒకటి}{!!{injective}} అని వ్యక్తం చేయడం కూడా ఇప్పుడు
కష్టం కాదు:
\begin{multline*}
 f \colon X \to Y \land \lforall[x][\lforall[x'][((x \in X \land x' \in
     X \land {}]] \\
 \lexists[y][(\mathrm{maps}(f, x, y) \land \mathrm{maps}(f,
        x', y))]) \lif x = x')
\end{multline*}
ప్రమేయం~$f\colon X\to Y$ \tetoken{ఒకటి-ఒకటి}{!!{injective}} అవుతుంది, అప్పుడు మాత్రమే:
$f$, $x,x'\in X$లను ఒకే~$y$కు చిత్రించినప్పుడల్లా $x=x'$ కావాలి. ఈ సూత్రాన్ని
$\mathrm{inj}(f,X,Y)$గా సంక్షిప్తీకరిస్తే, సమితుల సిద్ధాంత భాషలో
కాంటర్ సిద్ధాంతం వంటి అల్పం కాని విషయాన్ని ఇప్పటికే ప్రకటించగలం:
$\Pow{X}$ నుంచి $X$కు \tetoken{ఒకటి-ఒకటి}{!!{injective}} ప్రమేయం ఏదీ లేదు:
\[
\lforall[X][\lforall[Y][(\Pow{X} = Y \lif
    \lnot\lexists[f][\mathrm{inj}(f, Y, X)])]]
\]

ప్రతి నిర్వచక ధర్మానికీ ఒక సమితి ఉనికిని హామీ ఇచ్చే మరో స్వీకృతం
సమితుల సిద్ధాంతానికి అవసరమని అనిపించవచ్చు. $!A(x)$ అనేది స్వేచ్ఛా
చరం~$x$ గల సమితుల సిద్ధాంత సూత్రమైతే, కింది \tetoken{వాక్యాన్ని}{!!{sentence}}
పరిశీలించవచ్చు:
\[
\lexists[y][\lforall[x][(x \in y \liff !A(x))]].
\]
ఈ \tetoken{వాక్యం}{!!{sentence}}, సరిగ్గా $!A(x)$ను సంతృప్తిపరిచే~$x$లే \tetoken{మూలకాలుగా}{!!{element}s}
ఉన్న ఒక సమితి~$y$ ఉందని చెబుతుంది. ఈ పథకాన్ని ``ధర్మసంగ్రహ సూత్రం''
అంటారు. ఇది చాలా ఉపయోగకరంగా కనిపిస్తుంది; కానీ దురదృష్టవశాత్తూ ఇది
అవైరుధ్యమైనది కాదు. $!A(x) \ident \lnot x \in x$గా తీసుకుంటే,
ధర్మసంగ్రహ సూత్రం కింది వాక్యాన్ని ప్రకటిస్తుంది:
\[
\lexists[y][\lforall[x][(x \in y \liff x \notin x)]],
\]
అంటే తమలో తాము \tetoken{మూలకాలు}{!!{element}s} కాని అన్ని సమితుల సమితి ఉనికిని
ప్రకటిస్తుంది. అలాంటి సమితి ఉండజాలదు---ఇదే రసెల్ వైరుధ్యం.
నిజానికి, $\Log{ZFC}$లో ఈ సూత్రానికి పరిమితం చేసిన---మరియు అవైరుధ్యమైన---
రూపమైన వేరుచేయు సూత్రం ఉంది:
\[
\lforall[z][\lexists[y][\lforall[x][(x \in y \liff (x \in z \land
      !A(x))]]].
\]

\begin{prob}
కింది దాన్ని చూపే \tetoken{ఒక వ్యుత్పత్తిని}{!!a{derivation}} ఇవ్వడం ద్వారా ధర్మసంగ్రహ సూత్రం
అవైరుధ్యమైనది కాదని చూపండి:
\[
\lexists[y][\lforall[x][(x \in y \liff x \notin x)]] \Proves \lfalse.
\]
మొదట $(A \lif \lnot A) \land (\lnot A \lif A)
\Proves \lfalse$ అని చూపడం ఉపయోగపడవచ్చు.
\end{prob}
\end{document}
